\SourceLeaf{122}
\[
\begin{gathered}
M\bigl(1-e^{\rho\rho}\surd\tfrac\pi{8m}\bigr)\quad\text{und}\quad M\bigl(1+e^{\rho\rho}\surd\tfrac\pi{8m}\bigr),\\[4pt]
\text{oder in Zahlen}\quad M\bigl(1\mp\tfrac{0{,}7520974}{\surd m}\bigr)
\end{gathered}
\]
Dies Verfahren ist also nur wenig genauer, als die Anwendung der Formel VI, und man müsste 249 Beobachtungsfehler zu Rathe ziehen, um eben so weit zu reichen, wie mit 100 Beobachtungsfehlern nach Formel II.

\Article{8}
Die Anwendung einiger von diesen Methoden auf die in \textsc{Bode}'s astronomischem Jahrbuche für 1818 S. 234 vorkommenden Fehler bei 48 Beobachtungen der geraden Aufsteigungen des Polarsterns von \textsc{Bessel} gab
\[
S'=60''46,\quad S''=110''600,\quad S'''=250''341118
\]
Hieraus folgten die wahrscheinlichsten Werthe von \(r\)
\begin{center}
\begin{tabular}{r@{\;}c@{\;}l@{\quad}l@{\;}l}
nach Formel I&.\,.\,.\,.&\(1''065,\)&wahrscheinl. Unsicherheit \(=\)&\(\pm0''078\)\\
II&.\,.\,.\,.&\(1{,}024\)&&\(\pm0{,}070\)\\
III&.\,.\,.\,.&\(1{,}001\)&&\(\pm0{,}072\)\\
nach Art. 7&.\,.\,.\,.\,.&\(1{,}045\)&&\(\pm0{,}113\)
\end{tabular}
\end{center}
eine Uebereinstimmung, wie sie kaum zu erwarten war. \textsc{Bessel} giebt selbst \(1''067\), und scheint daher der Formel I gemäss gerechnet zu haben.
\begin{center}\rule{.24\linewidth}{.3pt}\end{center}
